\emph{Jason Mark Wrinkle (principal investigator and first author);
Yolanda Biche et al.~(drafting and computation). Working chapter,
2026-11-21. Status markers are kept visible: claims about the project's
own state are dated, and unsettled design questions are named as open
rather than smoothed over.}

\begin{center}\rule{0.5\linewidth}{0.5pt}\end{center}

\section{1.1 The claim this chapter
supports}\label{the-claim-this-chapter-supports}

TaxoPolemics is an engine for mapping disagreement. Given a domain of
difficult questions, it builds a \emph{hypermap}: a taxonomy of topics,
each topic articulated into two to eight \emph{poles} --- the stances
the topic admits --- and each person located in that space by rating,
per topic, how far they lean toward each pole and, independently, how
much the question matters to them. Two such profiles can then be
\emph{compared} by formal operations that distinguish two relations
ordinary language conflates: whether two people could \textbf{commune}
(the crossed register) and whether they could \textbf{converse} (the
straight register).

This chapter argues that the engine is not an orphan. It descends from
at least five identifiable traditions --- classical stasis theory, the
modern theory of argument, the design-methods tradition of issue-based
information systems, the study of classification itself, and the
psychometrics of attitude --- and that in each case the descent is
structural, not decorative: the ancestor solved a piece of the same
problem, and TaxoPolemics inherits both the solution and the problem it
left behind. A framework that can name its ancestors can also name its
delta. The final section does so.

\section{1.2 Stasis theory: the first taxonomy of
disagreement}\label{stasis-theory-the-first-taxonomy-of-disagreement}

The oldest ancestor is the most direct. In the second century BCE,
Hermagoras of Temnos founded what the rhetorical tradition calls
\textbf{stasis theory} (Greek \emph{stasis}, a stand, a position taken
in a dispute). Stasis theory holds that any controversy, to be argued at
all, must first be \emph{located}: the disputants must establish what
kind of question they are actually disagreeing about. The canonical
system distinguishes four stases --- \textbf{conjecture} (did it
happen?), \textbf{definition} (what is it?), \textbf{quality} (was it
right, how serious is it?), and \textbf{procedure} (is this the right
forum, the right process?). Quintilian and Cicero both transmit the
doctrine; it structured rhetorical education for over a millennium.

Stasis theory is, in modern vocabulary, a \emph{taxonomy of
disagreement} --- a claim that disputes have kinds, that the kinds can
be enumerated, and that identifying the kind is prior to, and more
valuable than, winning the argument. Every subsequent rhetorical
treatment of ``finding the real issue'' is downstream of it.

TaxoPolemics inherits the founding conviction --- that the first
intellectual act in a dispute is locative, not combative --- and changes
the unit of analysis. Stasis classifies \emph{disputes}; TaxoPolemics
classifies \emph{stances within topics}, and computes locations for
\emph{persons}. Where the stasis system asks ``what is this argument
about?'', the hypermap asks ``on this question, which poles exist, where
does each person stand among them, and how much does each person care?''
The ancient system is a diagnostic art exercised by a trained rhetor;
the modern one is an instrument any two disputants can run on
themselves. The continuity is the point: both rest on the premise that
most failed arguments fail because the parties never established what
they were disagreeing about, and that a shared map of the disagreement
is a precondition of a productive one.

\section{1.3 The modern theory of argument: Toulmin, Perelman,
Burke}\label{the-modern-theory-of-argument-toulmin-perelman-burke}

Three twentieth-century rhetoricians rebuilt the study of argument after
its long subordination to formal logic, and each contributes a
load-bearing element.

\textbf{Stephen Toulmin's} \emph{The Uses of Argument} (1958) proposed
that arguments have a \emph{layout} --- claim, grounds, warrant,
backing, qualifier, rebuttal --- and, more importantly for our purposes,
that the force of argumentative terms is \textbf{field-dependent}:
``probable,'' ``sufficient,'' ``reasonable'' mean different things in
law, in science, and at the dinner table, and no field-independent logic
governs them all. The hypermap system is field-dependent in exactly
Toulmin's sense, and wears the fact on its face: one engine runs sixteen
domains --- sexuality, politics, Abrahamic faith, chemistry, aesthetics,
philosophy, and others --- and in each domain the \emph{vocabulary of
the comparison itself} is re-skinned. The crossed relation is ``Play''
in the first domain and ``Commune'' in the vanilla domains; the straight
relation is ``Parlay'' in one and ``Converse'' in the others; the depth
measure is the Kink Factor in one and the Depth Factor elsewhere. A
field-independent framework would not need such skinning. TaxoPolemics
assumes Toulmin was right: the instrument must speak each field's own
language or it will not be believed in any of them.

\textbf{Chaïm Perelman and Lucie Olbrechts-Tyteca's} \emph{The New
Rhetoric} (1958) restored \emph{audience} to the centre of argument
theory: argumentation is the study of the means by which minds gain
adherence, and its basic data are not propositions but \emph{agreements}
--- the shared premises, values, and hierarchies from which a particular
audience can be moved. Their catalogue of the ``starting points of
argumentation'' is, like stasis, a taxonomy, and like this project, it
is interested less in who wins than in the prior question of what two
parties already hold in common. The straight register of the hypermap
--- the measurement of shared ground, in the same orientation, between
two profiles --- is Perelman's question rendered computable: not ``can
these two people persuade each other?'' but the antecedent ``what,
exactly, do they already agree on, and how strongly?''

\textbf{Kenneth Burke} contributes the project's central distinction. In
\emph{A Grammar of Motives} (1945) and the \emph{Rhetoric of Motives}
(1950), Burke argued that rhetoric is rooted in \textbf{identification}:
persuasion works not by defeating an opponent but by establishing
consubstantiality --- a sense of shared substance --- across division.
Identification is the correct ancestor for the crossed register's ruled
meaning on the vanilla maps: \textbf{communion}. The crossed comparison,
in which each person's leaning is read against the complementary leaning
of the other, asks whether two people can be \emph{of one substance} on
a topic --- can live it, work it, be it together --- as distinct from
the straight comparison, which asks only whether they would find each
other interesting to talk to. That the two relations are different
facts, priced separately by different mathematics (an indefinite form
for the crossed, a positive-definite one for the straight, as the
companion paper shows), is the project's central formal claim, and it is
a Burkean claim before it is a mathematical one.

\section{1.4 The design-methods tradition: IBIS and the wicked
problem}\label{the-design-methods-tradition-ibis-and-the-wicked-problem}

The second lineage is not rhetorical but computational, and it is the
closest technical ancestor. In the 1970s, Horst Rittel --- who had
already named the class of \textbf{wicked problems}, problems with no
definitive formulation and no stopping rule --- built, with Werner Kunz,
the \textbf{Issue-Based Information System (IBIS)}: a notation for
design deliberation with exactly three node types. An \textbf{Issue} is
a question under dispute; \textbf{Positions} are the stances one can
take on it; \textbf{Arguments} support or object to positions. IBIS was
implemented in software (notably Jeff Conklin's gIBIS in the 1980s) and
survives today in dialogue-mapping practice.

The resemblance to a hypermap topic is not accidental and should not be
understated: an IBIS Issue \emph{is} a topic, its Positions \emph{are}
poles, and the design intent --- externalise the structure of a
contested question so a group can deliberate over a shared map rather
than past each other --- is the same intent. TaxoPolemics is IBIS's
descendant, and the descent specifies the delta with unusual clarity:

\begin{enumerate}
\def\labelenumi{\arabic{enumi}.}
\tightlist
\item
  \textbf{Measurement.} IBIS positions are categorical: one adopts a
  position or one does not. A hypermap pole carries a signed magnitude
  --- a leaning, -100 to +100 --- so agreement and opposition become
  matters of degree, and two people who nominally hold ``the same
  position'' can be shown to hold it at different strengths and in
  different blends.
\item
  \textbf{The (n+1)th coordinate.} IBIS has no representation of
  \emph{how much an issue matters} to a participant, independent of the
  position taken on it. TaxoPolemics separates the leanings (n numbers
  for n poles) from \textbf{intensity} --- a signed relevance
  coordinate, the (n+1)th number --- as a matter of ruled design
  (2026-11-19). The separation is the project's single most
  consequential modelling decision, and \S 1.6 shows it has an independent
  ancestry in attitude research.
\item
  \textbf{Comparison as a first-class operation.} IBIS maps a
  deliberation; it does not compare two deliberators. The hypermap's
  pair mathematics --- the tensor product of profiles, its crossed and
  straight contractions, and the scores they yield --- has no IBIS
  analogue. IBIS tells a group what the questions are; TaxoPolemics
  additionally tells two people where they stand relative to each other,
  with receipts.
\item
  \textbf{Generativity of topics.} IBIS issues are authored in the room.
  Hypermap taxonomies are authored in advance, as \emph{skins} ---
  curated instruments, one per domain, sharing an engine. The
  deliberation is therefore instrumented rather than improvised: closer
  to a psychometric test than to a flip-chart, with the authoring burden
  (and its responsibilities, treated in Chapter Two) falling on the
  taxonomy's designer.
\end{enumerate}

\section{1.5 Knowledge organization: classification as
argument}\label{knowledge-organization-classification-as-argument}

The third lineage is the study of classification itself. Geoffrey Bowker
and Susan Leigh Star's \emph{Sorting Things Out: Classification and Its
Consequences} (1999) established, for information science, that
classification systems are neither neutral nor innocent: every taxonomy
encodes a theory of its domain, renders some distinctions visible and
others unthinkable, and has consequences for the people sorted by it.
The field of \textbf{knowledge organization} treats the design of a
classification --- its categories, its facets, its rules of combination
--- as scholarly work in its own right, judged by coherence, coverage,
and candour about its own choices.

A hypermap skin is a classification in exactly this sense, and the
project accepts the field's discipline. Two of its standing rules are
knowledge-organization rules. The \textbf{neutral umbrella rule} (ruled
2026-11-17) governs topics whose names join two stances with an
ampersand or slash: the joined name must be a \emph{neutral subject term
covering both sides}, never one of the joined terms itself --- a topic
may not smuggle a verdict into its own title. And the authoring standard
for topic descriptions (the \emph{hovertip}: the issue stated neutrally,
then one line per pole, each stated as its holders would state it) is a
rule of representational fairness of the kind Bowker and Star demand and
few production taxonomies observe. That a sexuality taxonomy, a politics
taxonomy, and a consciousness taxonomy (the sixteenth skin, Nouscapades,
added 2026-11-21) are built under one set of such rules is itself a
claim of the work: classification design is portable across domains
precisely when its fairness rules are stated at the level of the engine,
not improvised per domain.

\section{1.6 Attitude research: the independent ancestry of the (n+1)
model}\label{attitude-research-the-independent-ancestry-of-the-n1-model}

The separation of \emph{direction} from \emph{intensity} --- of which
way a person leans from how much the question matters to them --- was
arrived at in this project from instrument design (a slider's blend and
a slider's brightness are different facts, and conflating them corrupts
both). It is therefore worth recording that social psychology reached
the same separation from the opposite direction, empirically. The
\textbf{attitude strength} literature (Krosnick, Petty, and
collaborators, from the 1990s) established that an attitude's
\emph{direction and extremity} and its \emph{strength} --- importance,
certainty, accessibility --- are distinct properties with distinct
consequences: strong attitudes persist, resist persuasion, and guide
behaviour; weak ones, however extreme on paper, do not. A poll that
reports only leanings is therefore systematically misleading, in a way
the literature can quantify.

TaxoPolemics' (n+1) model is the attitude-strength distinction built
into the instrument rather than corrected for afterwards: every topic
rating is a vector of leanings \emph{plus} a signed intensity ---
positive for caring in favour of the issue's salience, negative for
active aversion to the issue itself --- displayed as brightness (full at
maximum positive, dimming to indifference at zero, passing into
complementary colours below zero). The companion tensor paper states the
pair mathematics for the two-flank case; the multi-pole generalisation,
including the open question of how the intensity coordinate should
weight pair scoring, is stated there as an open problem rather than
resolved by fiat. An honest lineage includes the note that the project's
most original modelling decision is also its least finished --- a state
of affairs Chapter Three's defence treats as a feature of the work's
stage, not an embarrassment to be hidden.

\section{1.7 Further kin: deliberative democracy and the map as
medium}\label{further-kin-deliberative-democracy-and-the-map-as-medium}

Two nearer relations deserve brief acknowledgement. \textbf{Deliberative
polling} (James Fishkin, from 1988) demonstrated that structured
exposure to balanced information changes what citizens think in
measurable, non-random ways --- evidence that disagreement is partly an
artefact of unstructured encounter, and that instruments can improve it.
The hypermap is a deliberative instrument in this tradition, though a
self-administered and pairwise one rather than a convened one. And the
project's commitment to the map as the primary object --- the ratings
are drawn as a navigable hypermap, not tabulated as a report ---
descends from the argument-mapping and information-visualisation
tradition (and from Tufte's insistence that a display is an argument), a
commitment Chapter Two treats methodologically: in this system, the
visualisation is not an illustration of the data but the form in which
the data is \emph{read}, and several of the project's defects (recorded
in the companion paper, \S 7) were precisely cases where a display quietly
stopped being the statistic it claimed to be.

\section{1.8 The delta, stated}\label{the-delta-stated}

Against these ancestors, the work claims four things as its own:

\begin{enumerate}
\def\labelenumi{\arabic{enumi}.}
\tightlist
\item
  \textbf{A measured stance-space.} Stasis located disputes; IBIS mapped
  them; neither measured the disputants. The hypermap locates
  \emph{persons} in a pole space with signed magnitudes, per topic, at
  scales from two to eight poles.
\item
  \textbf{Intensity as an independent coordinate.} The (n+1) model ---
  leanings plus a signed relevance --- as a ruled, first-class part of
  the rating, not a derived quantity.
\item
  \textbf{Two priced relations.} The crossed and straight comparisons
  --- communion and conversation --- as distinct mathematical objects
  (indefinite and positive-definite forms respectively) computed over
  the same product, where every predecessor system that scores
  ``compatibility'' at all produces a single number that confuses the
  two.
\item
  \textbf{Receipts.} Every headline number in the system can be
  \emph{un-contracted} --- opened back into the per-topic components,
  gates, and settings that produced it. A compatibility figure that
  cannot show its receipts is a horoscope; the project's founding
  discipline, stated in the companion paper and enforced throughout its
  development, is that no number ships without them.
\end{enumerate}

What the work does \emph{not} claim: that its taxonomies are complete,
that its scoring definitions are final (the six pair scores stand as
provisionally ruled definitions, 2026-10-30, pending review), or that
the instrument has been validated at population scale. Those are the
subjects of the next two chapters.

\begin{center}\rule{0.5\linewidth}{0.5pt}\end{center}

\emph{Next: Chapter Two --- Methodology: research through design,
taxonomy construction, the measurement model, and the evaluation
regime.}
